\chapter{Apéndice estadístico}

Fuente de todas las tablas: elaboración propia.

\section{Estadísticos descriptivos}

Muestra efectiva de estimación, 1998T2-2025T4 (111 trimestres).

{\def\LTcaptype{none} % do not increment counter
\begin{longtable}[]{@{}
  >{\raggedright\arraybackslash}p{(\linewidth - 8\tabcolsep) * \real{0.2000}}
  >{\raggedright\arraybackslash}p{(\linewidth - 8\tabcolsep) * \real{0.2000}}
  >{\raggedright\arraybackslash}p{(\linewidth - 8\tabcolsep) * \real{0.2000}}
  >{\raggedright\arraybackslash}p{(\linewidth - 8\tabcolsep) * \real{0.2000}}
  >{\raggedright\arraybackslash}p{(\linewidth - 8\tabcolsep) * \real{0.2000}}@{}}
\toprule\noalign{}
\begin{minipage}[b]{\linewidth}\raggedright
Variable
\end{minipage} & \begin{minipage}[b]{\linewidth}\raggedright
Media
\end{minipage} & \begin{minipage}[b]{\linewidth}\raggedright
Desvío estándar
\end{minipage} & \begin{minipage}[b]{\linewidth}\raggedright
Mínimo
\end{minipage} & \begin{minipage}[b]{\linewidth}\raggedright
Máximo
\end{minipage} \\
\midrule\noalign{}
\endhead
\bottomrule\noalign{}
\endlastfoot
Resultado primario (\% del PIB, trimestral anualizado) & -0,04 & 3,23 &
-12,16 & 7,02 \\
Crecimiento real trimestral (\%) & 0,39 & 7,41 & -13,77 & 18,21 \\
EMBI+ (puntos básicos) & 1.561,19 & 1.611,94 & 210,22 & 6.658,98 \\
TCRM (diciembre de 2001 = 1) & 1,66 & 0,45 & 0,98 & 2,76 \\
Deuda consolidada neta, rezagada (\% del PIB) & 57,74 & 27,07 & 26,13 &
136,60 \\
Variación de la tasa de la FED (puntos) & -0,01 & 0,52 & -1,43 & 1,69 \\
VIX & 20,20 & 7,33 & 10,31 & 58,60 \\
\end{longtable}
}

\section{Propiedades de las series}

\textbf{Raíz unitaria.} ADF (Dickey y Fuller, 1979; hipótesis nula: raíz
unitaria) y KPSS (Kwiatkowski et al., 1992; hipótesis nula:
estacionariedad), con constante. Una serie se clasifica como
estacionaria solo si ADF rechaza y KPSS no rechaza al 5\%, y con raíz
unitaria en el caso inverso. Zivot y Andrews (1992) permiten un quiebre
endógeno en el nivel. Muestra 1996T4-2025T4.

{\def\LTcaptype{none} % do not increment counter
\begin{longtable}[]{@{}
  >{\raggedright\arraybackslash}p{(\linewidth - 10\tabcolsep) * \real{0.1667}}
  >{\raggedright\arraybackslash}p{(\linewidth - 10\tabcolsep) * \real{0.1667}}
  >{\raggedright\arraybackslash}p{(\linewidth - 10\tabcolsep) * \real{0.1667}}
  >{\raggedright\arraybackslash}p{(\linewidth - 10\tabcolsep) * \real{0.1667}}
  >{\raggedright\arraybackslash}p{(\linewidth - 10\tabcolsep) * \real{0.1667}}
  >{\raggedright\arraybackslash}p{(\linewidth - 10\tabcolsep) * \real{0.1667}}@{}}
\toprule\noalign{}
\begin{minipage}[b]{\linewidth}\raggedright
Serie
\end{minipage} & \begin{minipage}[b]{\linewidth}\raggedright
ADF (p)
\end{minipage} & \begin{minipage}[b]{\linewidth}\raggedright
KPSS (p)
\end{minipage} & \begin{minipage}[b]{\linewidth}\raggedright
Clasificación
\end{minipage} & \begin{minipage}[b]{\linewidth}\raggedright
Zivot-Andrews (p)
\end{minipage} & \begin{minipage}[b]{\linewidth}\raggedright
Quiebre de Zivot-Andrews
\end{minipage} \\
\midrule\noalign{}
\endhead
\bottomrule\noalign{}
\endlastfoot
Resultado primario & 0,591 & 0,010 & raíz unitaria & 0,982 & 2011T2 \\
Crecimiento real & 0,053 & 0,100 & ambigua & 0,236 & 2003T1 \\
EMBI+ & 0,151 & 0,100 & ambigua & 0,087 & 2005T2 \\
TCRM & 0,250 & 0,100 & ambigua & 0,000 & 2001T4 \\
Deuda consolidada neta & 0,099 & 0,100 & ambigua & 0,575 & 2005T1 \\
Tasa de la FED (nivel) & 0,019 & 0,047 & ambigua & 0,020 & 2021T2 \\
VIX & 0,000 & 0,100 & estacionaria & 0,028 & 2019T4 \\
\end{longtable}
}

\textbf{Quiebres estructurales.} El procedimiento de Bai y Perron (2003)
sobre la media de la deuda consolidada neta, con el número de quiebres
elegido por el criterio de Schwarz, encuentra 3 quiebres: 2001T4, 2006T2
y 2018T1. Son los que entran como variables binarias en el SVARX.

\textbf{Cointegración.} Prueba de la traza de Johansen (1991) sobre las
variables en niveles (resultado primario, EMBI+, TCRM y deuda), con los
rezagos del VAR. El crecimiento se excluye porque es una tasa de
variación: incluido, genera por sí solo una relación de cointegración
trivial.

{\def\LTcaptype{none} % do not increment counter
\begin{longtable}[]{@{}llll@{}}
\toprule\noalign{}
Hipótesis nula & Traza & Valor crítico 5\% & Rechaza \\
\midrule\noalign{}
\endhead
\bottomrule\noalign{}
\endlastfoot
rango ≤ 0 & 43,7 & 47,9 & no \\
rango ≤ 1 & 20,0 & 29,8 & no \\
rango ≤ 2 & 10,8 & 15,5 & no \\
rango ≤ 3 & 3,3 & 3,8 & no \\
\end{longtable}
}

\section{Selección del orden del VAR}

Muestra efectiva común 1998T2-2025T4 (111 trimestres) para todos los
órdenes (sección 4.5). El orden elegido es el mayor orden estable que no
supera al que minimiza el AIC.

\textbf{Deuda consolidada neta, valor nominal (principal).} Orden
elegido: 3.

{\def\LTcaptype{none} % do not increment counter
\begin{longtable}[]{@{}lllllll@{}}
\toprule\noalign{}
p & AIC & BIC & HQ & Raíz máxima & Estable & Portmanteau (p) \\
\midrule\noalign{}
\endhead
\bottomrule\noalign{}
\endlastfoot
1 & 12,351 & 13,914 & 12,985 & 0,855 & sí & 0,000 \\
2 & 12,052 & 14,005 & 12,844 & 0,894 & sí & 0,039 \\
3 & 12,125 & 14,468 & 13,075 & 0,964 & sí & 0,036 \\
4 & 12,009 & 14,743 & 13,118 & 1,021 & no & 0,023 \\
5 & 11,939 & 15,064 & 13,207 & 1,024 & no & 0,002 \\
6 & 11,918 & 15,433 & 13,344 & 1,033 & no & 0,000 \\
\end{longtable}
}

\textbf{Deuda consolidada neta, valor libros.} Orden elegido: 3.

{\def\LTcaptype{none} % do not increment counter
\begin{longtable}[]{@{}lllllll@{}}
\toprule\noalign{}
p & AIC & BIC & HQ & Raíz máxima & Estable & Portmanteau (p) \\
\midrule\noalign{}
\endhead
\bottomrule\noalign{}
\endlastfoot
1 & 12,326 & 13,888 & 12,960 & 0,855 & sí & 0,000 \\
2 & 12,043 & 13,996 & 12,835 & 0,900 & sí & 0,034 \\
3 & 12,122 & 14,466 & 13,073 & 0,961 & sí & 0,030 \\
4 & 12,003 & 14,737 & 13,112 & 1,014 & no & 0,021 \\
5 & 11,937 & 15,062 & 13,205 & 1,017 & no & 0,002 \\
6 & 11,918 & 15,433 & 13,344 & 1,015 & no & 0,000 \\
\end{longtable}
}

\textbf{Deuda consolidada bruta.} Orden elegido: 4.

{\def\LTcaptype{none} % do not increment counter
\begin{longtable}[]{@{}lllllll@{}}
\toprule\noalign{}
p & AIC & BIC & HQ & Raíz máxima & Estable & Portmanteau (p) \\
\midrule\noalign{}
\endhead
\bottomrule\noalign{}
\endlastfoot
1 & 12,378 & 13,940 & 13,011 & 0,818 & sí & 0,000 \\
2 & 12,082 & 14,034 & 12,874 & 0,840 & sí & 0,071 \\
3 & 12,170 & 14,514 & 13,121 & 0,923 & sí & 0,035 \\
4 & 12,057 & 14,791 & 13,167 & 0,987 & sí & 0,020 \\
5 & 11,984 & 15,108 & 13,251 & 1,004 & no & 0,002 \\
6 & 11,968 & 15,483 & 13,394 & 1,027 & no & 0,000 \\
\end{longtable}
}

\section{Diagnósticos de los residuos}

Portmanteau con 12 rezagos (autocorrelación), Jarque-Bera multivariada
(normalidad), ARCH-LM con cuatro rezagos por ecuación
(heterocedasticidad condicional) y estabilidad (1 = estable).

{\def\LTcaptype{none} % do not increment counter
\begin{longtable}[]{@{}
  >{\raggedright\arraybackslash}p{(\linewidth - 6\tabcolsep) * \real{0.2500}}
  >{\raggedright\arraybackslash}p{(\linewidth - 6\tabcolsep) * \real{0.2500}}
  >{\raggedright\arraybackslash}p{(\linewidth - 6\tabcolsep) * \real{0.2500}}
  >{\raggedright\arraybackslash}p{(\linewidth - 6\tabcolsep) * \real{0.2500}}@{}}
\toprule\noalign{}
\begin{minipage}[b]{\linewidth}\raggedright
Prueba (p-valor salvo estabilidad)
\end{minipage} & \begin{minipage}[b]{\linewidth}\raggedright
neta, valor nominal (principal)
\end{minipage} & \begin{minipage}[b]{\linewidth}\raggedright
neta, valor libros
\end{minipage} & \begin{minipage}[b]{\linewidth}\raggedright
bruta
\end{minipage} \\
\midrule\noalign{}
\endhead
\bottomrule\noalign{}
\endlastfoot
Portmanteau (12 rezagos) & 0,036 & 0,030 & 0,020 \\
Normalidad multivariada (p-valor) & 0,000 & 0,000 & 0,000 \\
Estable (1 = sí) & 1,000 & 1,000 & 1,000 \\
ARCH-LM(4) resultado primario (p-valor) & 0,991 & 0,988 & 0,993 \\
ARCH-LM(4) crecimiento (p-valor) & 0,989 & 0,990 & 0,978 \\
ARCH-LM(4) EMBI (p-valor) & 0,998 & 0,998 & 0,999 \\
ARCH-LM(4) TCRM (p-valor) & 0,867 & 0,824 & 0,395 \\
\end{longtable}
}

\section{Matriz de impacto estructural}

Efecto en el mismo trimestre de un \emph{shock} estructural de un desvío
estándar (especificación principal). Los ceros son las restricciones de
identificación de la sección 4.4. Los dos \emph{shocks} del bloque
cambiario y de riesgo no tienen interpretación por separado: solo se usa
su efecto conjunto.

{\def\LTcaptype{none} % do not increment counter
\begin{longtable}[]{@{}
  >{\raggedright\arraybackslash}p{(\linewidth - 8\tabcolsep) * \real{0.2000}}
  >{\raggedright\arraybackslash}p{(\linewidth - 8\tabcolsep) * \real{0.2000}}
  >{\raggedright\arraybackslash}p{(\linewidth - 8\tabcolsep) * \real{0.2000}}
  >{\raggedright\arraybackslash}p{(\linewidth - 8\tabcolsep) * \real{0.2000}}
  >{\raggedright\arraybackslash}p{(\linewidth - 8\tabcolsep) * \real{0.2000}}@{}}
\toprule\noalign{}
\begin{minipage}[b]{\linewidth}\raggedright
Variable
\end{minipage} & \begin{minipage}[b]{\linewidth}\raggedright
\emph{Shock} fiscal
\end{minipage} & \begin{minipage}[b]{\linewidth}\raggedright
\emph{Shock} de actividad
\end{minipage} & \begin{minipage}[b]{\linewidth}\raggedright
\emph{Shock} cambiario y de riesgo 1
\end{minipage} & \begin{minipage}[b]{\linewidth}\raggedright
\emph{Shock} cambiario y de riesgo 2
\end{minipage} \\
\midrule\noalign{}
\endhead
\bottomrule\noalign{}
\endlastfoot
Resultado primario & 1,509 & 0,679 & 0,000 & 0,000 \\
Crecimiento & -0,266 & 2,717 & 0,000 & 0,000 \\
EMBI+ & -118,560 & -30,491 & 449,797 & 0,000 \\
TCRM & 0,022 & -0,014 & 0,011 & 0,096 \\
\end{longtable}
}

\section{Ecuación del resultado primario}

Ecuación del resultado primario del VARX en forma reducida,
especificación principal (MCO, 1998T2-2025T4). No se muestran la
constante, las variables binarias estacionales ni las de los quiebres.
Los errores estándar y los p-valores son los de MCO y se informan como
referencia: con regresores que pueden tener raíz unitaria su
distribución no es la estándar (Sims et al., 1990), por lo que la
inferencia del texto usa los intervalos del \emph{bootstrap}. El
coeficiente de la deuda rezagada es un efecto de impacto sobre el
resultado primario trimestral. Con los rezagos propios del resultado
primario (suma 0,729), el multiplicador de largo plazo de la estimación
por MCO es 0,051/(1 − 0,729) ≈ 0,19; no se calculó su intervalo.

{\def\LTcaptype{none} % do not increment counter
\begin{longtable}[]{@{}
  >{\raggedright\arraybackslash}p{(\linewidth - 6\tabcolsep) * \real{0.2500}}
  >{\raggedright\arraybackslash}p{(\linewidth - 6\tabcolsep) * \real{0.2500}}
  >{\raggedright\arraybackslash}p{(\linewidth - 6\tabcolsep) * \real{0.2500}}
  >{\raggedright\arraybackslash}p{(\linewidth - 6\tabcolsep) * \real{0.2500}}@{}}
\toprule\noalign{}
\begin{minipage}[b]{\linewidth}\raggedright
Regresor
\end{minipage} & \begin{minipage}[b]{\linewidth}\raggedright
Coeficiente
\end{minipage} & \begin{minipage}[b]{\linewidth}\raggedright
Error estándar
\end{minipage} & \begin{minipage}[b]{\linewidth}\raggedright
p-valor
\end{minipage} \\
\midrule\noalign{}
\endhead
\bottomrule\noalign{}
\endlastfoot
Variación de la tasa de la FED, rezago 0 & 0,5297 & 0,5178 & 0,306 \\
Variación de la tasa de la FED, rezago 1 & -0,7911 & 0,5670 & 0,163 \\
Variación de la tasa de la FED, rezago 2 & 0,0030 & 0,5160 & 0,995 \\
VIX & -0,0253 & 0,0319 & 0,427 \\
Deuda consolidada neta, rezago 1 (reacción fiscal de impacto) & 0,0505 &
0,0215 & 0,019 \\
Resultado primario, rezago 1 & 0,3404 & 0,1116 & 0,002 \\
Crecimiento, rezago 1 & 0,0901 & 0,0647 & 0,164 \\
EMBI+, rezago 1 & -0,0005 & 0,0004 & 0,189 \\
TCRM, rezago 1 & 3,6341 & 1,7183 & 0,034 \\
Resultado primario, rezago 2 & 0,3739 & 0,1190 & 0,002 \\
Crecimiento, rezago 2 & 0,1250 & 0,0633 & 0,048 \\
EMBI+, rezago 2 & 0,0001 & 0,0005 & 0,850 \\
TCRM, rezago 2 & -2,9019 & 2,3871 & 0,224 \\
Resultado primario, rezago 3 & 0,0149 & 0,1154 & 0,897 \\
Crecimiento, rezago 3 & -0,0330 & 0,0605 & 0,586 \\
EMBI+, rezago 3 & 0,0002 & 0,0004 & 0,520 \\
TCRM, rezago 3 & -0,2936 & 1,6455 & 0,858 \\
\end{longtable}
}

\section{Resultados completos de la especificación principal}

Estimación central con parámetros corregidos por sesgo (Kilian, 1998) y
20.000 trayectorias; intervalos de 200 réplicas con 1.000 trayectorias
cada una. Autovalor máximo del VAR: 0,964 estimado y 0,9998 corregido.
Las réplicas usan una corrección del sesgo más reducida que la
estimación central, de modo que los intervalos percentiles no quedan
centrados en ella y la estimación central se ubica cerca de su límite
superior (sección 4.5).

{\def\LTcaptype{none} % do not increment counter
\begin{longtable}[]{@{}
  >{\raggedright\arraybackslash}p{(\linewidth - 4\tabcolsep) * \real{0.3333}}
  >{\raggedright\arraybackslash}p{(\linewidth - 4\tabcolsep) * \real{0.3333}}
  >{\raggedright\arraybackslash}p{(\linewidth - 4\tabcolsep) * \real{0.3333}}@{}}
\toprule\noalign{}
\begin{minipage}[b]{\linewidth}\raggedright
Medida
\end{minipage} & \begin{minipage}[b]{\linewidth}\raggedright
Estimación central
\end{minipage} & \begin{minipage}[b]{\linewidth}\raggedright
Intervalo de 90\%
\end{minipage} \\
\midrule\noalign{}
\endhead
\bottomrule\noalign{}
\endlastfoot
No estabilización 2031 (\%) & 25,2 & {[}19,0; 27,7{]} \\
No estabilización 2035 (\%) & 22,3 & {[}15,6; 24,9{]} \\
No estabilización 2031 con el resultado primario simulado (\%) & 29,2 &
{[}21,7; 31,1{]} \\
Probabilidad de que la deuda de 2031 supere a la de 2029 (\%) & 22,0 &
{[}13,2; 23,8{]} \\
Deuda mediana 2031 (\% del PIB) & 39,5 & {[}38,8; 40,3{]} \\
Deuda media 2035 (\% del PIB) & 31,3 & {[}29,5; 33,6{]} \\
Ancho del abanico 2031 (p5-p95) & 35,3 & {[}19,1; 37,6{]} \\
Varianza 2035: fiscal (\%) & 32,1 & {[}4,4; 52,4{]} \\
Varianza 2035: actividad (\%) & 9,3 & {[}3,8; 43,1{]} \\
Varianza 2035: cambiario y de riesgo (\%) & 58,7 & {[}17,2; 86,0{]} \\
No fiscal menos fiscal (puntos) & 35,9 & {[}-4,8; 91,2{]} \\
Cambiario y de riesgo menos fiscal (puntos) & 26,6 & {[}-31,2;
80,3{]} \\
Promedio por \emph{shock} cambiario y de riesgo menos fiscal (puntos) &
-2,5 & {[}-40,7; 38,1{]} \\
FED, nivel 0,50 pp más alto: diferencia en la probabilidad 2035 (pp) &
1,5 & {[}0,1; 3,1{]} \\
FED, nivel 0,50 pp más alto: diferencia en la deuda 2035 (pp del PIB) &
1,0 & {[}-0,1; 2,4{]} \\
FED, senda 2022-2025: diferencia en la probabilidad 2031 (pp) & 2,4 &
{[}-4,5; 11,4{]} \\
FED, senda 2022-2025: diferencia en la probabilidad 2035 (pp) & 13,0 &
{[}1,0; 24,8{]} \\
FED, senda 2022-2025: diferencia en la deuda 2035 (pp del PIB) & 7,8 &
{[}-1,4; 20,4{]} \\
\end{longtable}
}
